\chapter*{Legacy Problem Recovery Audit for VI/06}
\addcontentsline{toc}{chapter}{Legacy Problem Recovery Audit for VI/06}

The migration table assigns the unmatched remainder of the legacy CA13 source to VI/06 as a safety-net destination.  A
section-level reconstruction is necessary because that legacy file mixes the conceptual
prime-as-point bridge with material that belongs to later chapters.  The governing rule is
that unique mathematics is retained exactly once in its strongest canonical home.

\begin{longtable}{>{\raggedright\arraybackslash}p{0.21\textwidth}>{\raggedright\arraybackslash}p{0.34\textwidth}>{\raggedright\arraybackslash}p{0.31\textwidth}}
\toprule
\textbf{Legacy block} & \textbf{Topic} & \textbf{Canonical disposition}\\
\midrule
CA13 opening point discussion & Prime ideals as geometric points; ordinary versus generic behavior & VI/06 theory; retained and expanded\\
CA13 ``Can a prime ideal contain another ideal?'' & Ideal containment as incidence & VI/06.P1; retained and expanded\\
CA13 ``What is the ideal $(0)$ in $k[x]$?'' & Zero prime, prime but not maximal, generic-line intuition & VI/06.P2; retained and expanded\\
CA13 examples in $k[t]$, $\mathbb Z$, $k[x,y]/(xy)$ & Prime/maximal/minimal-prime intuition & VI/06 theory/examples/problems; retained selectively without duplicating later topology\\
CA13 Spec/Zariski sections & Definition and topology of $\Spec A$, closed sets, point closures & VI/07 and VI/09; reserved there\\
CA13 basic opens & $D(f)$ and basis of opens & VI/08; reserved there\\
CA13 local rings and residue fields & $A_{\mathfrak p}$ and $\kappa(\mathfrak p)$ & VI/11; reserved there\\
CA13 reduced/nonreduced examples & Nilpotents and doubled points & VI/10; reserved there\\
CA13 affine scheme/structure sheaf material & Scheme structure and regular functions & VI/18; reserved there\\
CA13 tensor/base-change section & Extension of scalars and splitting of primes & VI/24; one motivating example retained in VI/06.P6/C3, systematic theory reserved\\
CA13 quotient examples & Primes above an ideal and closed quotients & VI/06 quotient correspondence; scheme-theoretic closed subschemes reserved for VI/22\\
\bottomrule
\end{longtable}

\bigskip
\noindent\textbf{New canonical material.}
VI/06.P3--P8 and VI/06.C1--C2 are new synthesis items.  VI/06.C3 expands a legacy
field-extension warning into a controlled preview.  These additions do not replace legacy
material; they connect the classical affine chapters to the prime-spectrum architecture.

\bigskip
\noindent\textbf{Boundary rule.}
VI/06 answers \emph{why prime ideals deserve to be points}.  VI/07 answers \emph{how all
prime ideals become the topological space $\Spec A$}.  VI/08--VI/11 then develop its basic
opens, generic/closed points, nonreduced structure, local rings and residue fields.
